% The hwde showcase, in the pdf-material-builder house style (technical-doc,
% v5.1). Build with ./build_docs.sh pdf, which runs the skill's
% scripts/build.sh (lualatex, three passes). The flow charts in diagrams/ are
% built separately and placed here as they are.
\documentclass[11pt]{article}
\newcommand\hspaper{a4paper}
\usepackage{housestyle}
\usepackage{tabularx}

\hsslug{hwde showcase}
\fancyfoot[L]{\hsrunlabel{\textcopyright{} 2026 Ihsan Salari. MIT License.}}
\renewcommand{\sectionmark}[1]{\markboth{#1}{}}
\hssection{\leftmark}
\titleformat{\subsection}{\hssubheadfont\fontsize{13pt}{17pt}\selectfont\color{ink}}{}{0pt}{}
\titlespacing*{\subsection}{0pt}{16pt}{4pt}
\newcolumntype{Y}{>{\raggedright\arraybackslash}X}

\graphicspath{{renders/}}

% One board as a figure plate: file, name, grey facts, status, caption.
\newcommand{\board}[5]{\hsplate{#1}{#2}{#3}{#4}{#5}}
% Two half-measure plates side by side, sharing a baseline. At half measure
% the name and facts line is set ragged and never hyphenated.
\newcommand{\pairtext}{\hsfull\hyphenpenalty=10000\exhyphenpenalty=10000\relax}
\newcommand{\boardpair}[2]{%
  \begin{hsblock}%
  \begin{minipage}[t]{0.47\linewidth}\pairtext#1\end{minipage}\hfill
  \begin{minipage}[t]{0.47\linewidth}\pairtext#2\end{minipage}%
  \end{hsblock}}
\newcommand{\sep}{\ \textperiodcentered\ }
\newcommand{\renderfrom}[1]{Rendered by docs/showcase/render\_boards.sh (KiCad
  10.0.5) from the board file in boards/<board>/kicad, for #1. Layers, size and
  footprints from the same file, the stage from state.json, 23 September 2026.}

\begin{document}

\hstitleblock{An AI PCB engineer for KiCad and JLCPCB}%
{hwde}%
{Printed circuit boards, designed from a written brief.}

You describe the board in plain words. hwde researches the parts, picks a
topology, draws the schematic, lays out and routes the board, runs the
electrical and manufacturing checks, and hands you the files a factory needs.
A human engineer signs off at every stage.

\board{lumina-carrier}{lumina-carrier}%
{4 layers\sep 100 $\times$ 80 mm\sep 116 footprints}%
{ordered, fabricated}%
{A carrier board for a stage-lighting fixture that hwde designed.}

\hsprovenance{Rendered by docs/showcase/render\_boards.sh from
boards/lumina-carrier/kicad, 23 September 2026.}
\newpage

% ======================================================================
\section{How it works}
\hslead{Three pictures follow. The first is the route a board takes, the
second is who does which part of the work, and the third is what happens at
every handover.}

\begin{hsfigure}{Figure 1 / the route a board takes}%
{Each box in the row is a stage, and the box below it is the knowledge base
that the dashed arrows read. The orange dots mark where a person has to look at
the work and agree before it goes on.}
\hsdiagram{diagrams/pipeline}
\end{hsfigure}

A run starts from a written brief and ends with a manufacturing package of
Gerbers, a bill of materials and a pick-and-place file. Each stage in between
produces one file, and before it does, it reads the knowledge base of design
rules, datasheet numbers and what the factory can really make.

\subsection{Who does what}

The work is split in two. Twenty-two specialist subagents make the judgement
calls, like choosing a topology, sizing a part or reviewing a schematic, and
each one has a written contract for what it may decide. About sixty scripts do
and check the mechanical work, and they are all deterministic. They drive KiCad,
run the routers and export the Gerbers. A script never has an opinion, and an
agent never edits a board file by hand. Whatever either of them produces lands
as a file in one folder per board.

\hsprovenance{Figure 1: docs/showcase/diagrams/pipeline.tex. Subagents and
scripts: .claude/skills/hwde/agents and .claude/skills/hwde/scripts.}
\newpage

\begin{hsfigure}{Figure 2 / who does what}%
{The orange box is you and the blue one routes the work. Purple makes the
judgement calls, green does the mechanical work, and both write to the board
folder.}
\hsdiagram[width=0.72\linewidth]{diagrams/roles}
\end{hsfigure}

\subsection{The gate loop}

\begin{hsfigure}{Figure 3 / what happens at every handover}%
{The green box is the checks and the orange one is a person's sign-off. The
line from the checks down to the fix loop is the fail path, and it returns to
the file the stage made.}
\hsdiagram{diagrams/gate}
\end{hsfigure}

Every stage ends at a gate. The gate runs the checks that apply to that stage,
which are the electrical rules, the design rules, manufacturability and circuit
simulations against stated bounds. It either passes, or it hands the findings
back to a fix loop that reworks the file and tries again. A pass is not the end
of it, because a person still signs off before the next stage starts.

\hsprovenance{Figures 2 and 3: docs/showcase/diagrams/roles.tex and gate.tex.
The checks are the scripts in .claude/skills/hwde/scripts.}
\newpage

\subsection{How it gets better}

\begin{hsfigure}{Figure 4 / how one run feeds the next}%
{Solid arrows are one board run, from what broke to the notes it leaves. The
dashed line back is the next board starting from those notes.}
\hsdiagram{diagrams/lessons}
\end{hsfigure}

Every run writes down what broke and why, with a date and a tag, and that goes
into the knowledge base the next run reads first. The rule is to record a
gotcha the moment you hit it, so the next board does not pay for it again.

\hsclaim{hwde takes a board from a written brief to the files a factory needs,
and a person signs off at every stage.}{What hwde is for}

\begin{hscallout}{What a pass does not mean}
A green gate only proves the checks that actually ran, so the reports are kept
one per check instead of being collapsed into a verdict. The stage a board
reached is not a release certificate either. A person decides what is fit to
release.
\end{hscallout}

\hsprovenance{Figure 4: docs/showcase/diagrams/lessons.tex. The notes are
LEARNINGS.md.}
\newpage

% ======================================================================
\section{The boards}
\hslead{Fourteen boards got far enough to have a laid-out, routed PCB. They run
from a linear regulator with five parts to a 200\,W RF power stage.}

\begin{hsblock}
\hslabel{Table 1 / the fourteen boards and where each one stopped}\par\vspace{8pt}
\begin{tabularx}{\linewidth}{@{}l Y r r l@{}}
\hstoprule
\hshead{Board} & \hshead{What it is} & \hshead{Layers} & \hshead{Size (mm)} & \hshead{Stage reached}\\
\hstoprule
lumina-carrier & lighting carrier, PoE-powered & 4 & 100 $\times$ 80 & ordered, fabricated\\\hline
pd-trigger     & USB-C PD bench trigger        & 2 & 48 $\times$ 30   & ordered, fabricated\\\hline
lumina-par     & RGBW PAR daughter board       & 4 & 100 $\times$ 80 & verification\\\hline
rf-de-20m      & 20 MHz Class E GaN stage      & 4 & 120 $\times$ 80 & manufacturability\\\hline
sbuck-5v3a     & 5 V / 3 A synchronous buck    & 4 & 50 $\times$ 40   & package ready\\\hline
usb-buck       & USB device dev board          & 4 & 50 $\times$ 40   & package ready\\\hline
g0-sense       & USB-C sensor node             & 2 & 35.8 $\times$ 28.3 & package ready\\\hline
stm32-blinky   & minimal STM32 dev board       & 2 & 50 $\times$ 40   & package ready\\\hline
rf-term-150w   & 150 W RF dummy load head      & 2 & 26 $\times$ 20   & manufacturability\\\hline
bb-adc         & single-channel ADC            & 2 & 54.8 $\times$ 34.9 & manufacturability\\\hline
bb-amp         & sensor amplifier chain        & 2 & 48 $\times$ 28.3 & manufacturability\\\hline
bb-buck        & bare buck converter           & 2 & 35 $\times$ 25   & manufacturability\\\hline
bb-ldo         & bare linear regulator         & 2 & 34.7 $\times$ 34.7 & manufacturability\\\hline
bb-mcu         & bare microcontroller board    & 2 & 34.8 $\times$ 22.3 & manufacturability\\
\hstoprule
\end{tabularx}
\end{hsblock}

``Stage reached'' says where the board stopped in the pipeline, and it does not
say the board is fit to release. Two boards say ``ordered, fabricated'' because
their own files record a JLCPCB order number that tracked through to shipped.
The rest stopped with a checked manufacturing package in hand. Each size is the
bounding box of the board outline, and not every outline is a rectangle.

\hsprovenance{Read 23 September 2026 from boards/<board>/kicad/<board>.kicad\_pcb
(layers, size) and boards/<board>/state.json (stage).}
\newpage

% ======================================================================
\subsection{Designed, ordered and fabricated}

Every render below was made from the board file in the repository, by the
script that sits next to this document. lumina-carrier, on the first page, is
the universal carrier for a stage-lighting fixture. It takes power over
Ethernet through a PD controller and a 100\,V buck, and it carries a frozen
expansion connector that a fixture-specific daughter board plugs into. It is
the biggest board here, with 4 layers and 116 footprints, and another board
depends on the connector it sets. The other board that was ordered is a bench
tool.

\board{pd-trigger}{pd-trigger}%
{2 layers\sep 48 $\times$ 30 mm\sep 30 footprints}%
{ordered from JLCPCB, tracked to shipped}%
{A USB-C Power Delivery sink controller negotiates with a charger for 5, 9, 12,
15 or 20 volts, and you pick the profile on the board. The negotiated rail
comes out on a screw terminal rated for 5\,A, so up to 100\,W. The board
converts nothing. It just passes the rail through.}

\hsprovenance{\renderfrom{pd-trigger}}
\newpage

\subsection{Power}

\boardpair
{\board{sbuck-5v3a}{sbuck-5v3a}%
{4 layers\sep 50 $\times$ 40 mm\sep 44 footprints}{package ready}%
{A synchronous buck that takes 7--18\,V in and holds 5.0\,V $\pm$2\% at 3\,A
out. It is an open-frame power module.}}
{\board{rf-term-150w}{rf-term-150w}%
{2 layers\sep 26 $\times$ 20 mm\sep 6 footprints}{manufacturability}%
{A 50-ohm, 150\,W dummy load head for DC to 25\,MHz. The element bolts to a
heatsink the user supplies, so the board is only the RF launch and the
mechanical interface. It renders close to bare because none of its six
footprints has a 3D model the renderer could find, so the connector and the
element are not drawn.}}

\hsprovenance{\renderfrom{sbuck-5v3a and rf-term-150w}}
\newpage

% ----------------------------------------------------------------------
\subsection{Harder boards}

\board{rf-de-20m}{rf-de-20m}%
{4 layers\sep 120 $\times$ 80 mm\sep 80 footprints}%
{stopped at manufacturability checks}%
{A 20\,MHz Class E RF power stage that puts 200\,W into 50 ohms from a 40\,V
bus. A PWM drive arrives on an SMA input, a GaN gate driver buffers it, and one
ground-referenced eGaN FET switches. A resonant tank and an L-match transform
the 4.6-ohm load line to 50 ohms at the output SMA. The owner fixed the topology
and the two core parts after a costed trade study. hwde did the rest.}

\hsprovenance{\renderfrom{rf-de-20m}}
\newpage

\label{pg:par}
\board{lumina-par}{lumina-par}%
{4 layers\sep 100 $\times$ 80 mm\sep 158 footprints}%
{stopped at verification}%
{The RGBW PAR daughter board that stacks on the carrier. The carrier's
connector document is frozen, so this board takes it as given and may not
redefine it. Its power budget of 8.6 to 20\,W was re-derived here, and depends
on the power-over-Ethernet class the fixture gets. It is the densest board here
at 158 footprints, and the only one whose outline is not a rectangle. The
silkscreen warns that the board floats at the power-over-Ethernet potential,
because an earthed probe breaks the fixture's power negotiation.}

\hsprovenance{\renderfrom{lumina-par}}
\newpage

% ----------------------------------------------------------------------
\subsection{Microcontroller boards}

\board{g0-sense}{g0-sense}%
{2 layers\sep 35.8 $\times$ 28.3 mm\sep 30 footprints}%
{package ready, release attestation built}%
{A temperature and humidity node that runs off USB-C. USB-C supplies 5\,V and
an LDO makes 3.3\,V from it, and an STM32G030 on its internal oscillator reads
a Sensirion SHT4x over I\textsuperscript{2}C. The same bus comes out on a Qwiic
connector, the readings go out over a UART header, and you program it over
SWD. hwde designed this one end to end in a single unattended container run.}

\hsprovenance{\renderfrom{g0-sense}}
\newpage

\boardpair
{\board{usb-buck}{usb-buck}%
{4 layers\sep 50 $\times$ 40 mm\sep 28 footprints}{package ready}%
{An STM32F103 board that talks USB at full speed. The Micro-B socket carries
both power and data, and there is an LED, a button and an SWD header.}}
{\board{stm32-blinky}{stm32-blinky}%
{2 layers\sep 50 $\times$ 40 mm\sep 20 footprints}{package ready}%
{The simplest useful board here. It has an STM32F103, an 8\,MHz crystal, one
LED, an LDO and an SWD header.}}

\subsection{Building blocks}

These five boards are deliberately bare, each one carrying a single block and
nothing else. They are there to be studied and measured on a bench, so they
leave out protection, filtering, indicators and spare rails on purpose. All
five stopped at manufacturability checks.

\hsprovenance{\renderfrom{usb-buck and stm32-blinky}}
\newpage

\boardpair
{\board{bb-buck}{bb-buck}%
{2 layers\sep 35 $\times$ 25 mm\sep 20 footprints}{}%
{It takes 18--30\,V in and gives 5\,V at 2\,A out on a screw terminal. It
carries the converter and exactly what its datasheet requires.}}
{\board{bb-ldo}{bb-ldo}%
{2 layers\sep 34.7 $\times$ 34.7 mm\sep 5 footprints}{}%
{It takes 5\,V in and gives 3.3\,V at 500\,mA out, holding regulation in still
air on the board's own copper.}}
\boardpair
{\board{bb-mcu}{bb-mcu}%
{2 layers\sep 34.8 $\times$ 22.3 mm\sep 14 footprints}{}%
{One MCU with a 3.3\,V rail coming in, its debug header and four GPIO. There is
no regulator and no second rail.}}
{\board{bb-adc}{bb-adc}%
{2 layers\sep 54.8 $\times$ 34.9 mm\sep 25 footprints}{}%
{0--5\,V in on a screw terminal, and 12 bits or better at up to 10\,kSa/s out
over a 0.1-inch header. It trades speed for DC accuracy.}}

\hsprovenance{\renderfrom{bb-buck, bb-ldo, bb-mcu and bb-adc}}
\newpage

\boardpair
{\board{bb-amp}{bb-amp}%
{2 layers\sep 48 $\times$ 28.3 mm\sep 14 footprints}{}%
{It takes a 0--20\,mV differential bridge signal and gives 0--3.3\,V
single-ended out, from DC to 1\,kHz. The nominal gain is 165 V/V.}}
{\vspace{14pt}\hsfull\color{inkseventy}Two more boards stopped before layout
and have no PCB to show, a 5\,V/3\,A buck and a strobe daughter board. Both are
in the repository at the stage they reached.}

\boardpair
{\board{lumina-carrier-bottom}{lumina-carrier, underneath}{bottom view}{}%
{Every part sits on the top side of this board, so the back carries routing and
vias and nothing else. This is how the factory sees it.}}
{\board{lumina-par-flat}{lumina-par, straight down}{top-down view}{}%
{The daughter board on page~\pageref{pg:par}, from straight above, so the
routing shows. It has four layers and 158 footprints, routed to the same rules
as everything else.}}

\hsprovenance{\renderfrom{bb-amp, lumina-carrier and lumina-par}}
\newpage

% ======================================================================
\section{What it can and cannot do yet}
\hslead{hwde is an engineering assistant, and it does not release a board by
itself. A human engineer drives it and signs off at every stage.}

\begin{hsstatrow}[4]
\hsstat{14}{boards with a routed PCB}
\hsstat{2}{ordered from JLCPCB and fabricated}
\hsstat{200\,W}{out of the RF stage, into 50 ohms}
\hsstat{158}{footprints on lumina-par, the densest}
\end{hsstatrow}

\subsection{What holds}

\begin{itemize}[leftmargin=12pt,itemsep=2pt]
\item Fourteen boards went from a written brief to a routed layout, including a
      4-layer 200\,W RF stage and a novel 100\,W USB-C PD design.
\item Two of them were designed with it, ordered from a factory and fabricated,
      and their own files record the order numbers tracked through to shipped.
\item Every check is a script, and they all follow the same contract. The
      findings are kept per check rather than as a single verdict.
\item Everything a run produces is a file in that board's folder, and you can
      read it without hwde.
\end{itemize}

\subsection{What does not}

\begin{itemize}[leftmargin=12pt,itemsep=2pt]
\item The factory has no public API for its manufacturability review, so someone
      still does that step by hand in a browser.
\item Spending money is deliberately hard. One code path can place an order, and
      it handles 2-layer boards only. It refuses a board whose files already
      record an order, and it will not run without a fresh quote and a matching
      design hash. A person also has to type out a confirmation naming the
      board, the quantity and the total.
\item The electrical and thermal checks are engineering screens, and each states
      its own accuracy. None of them is a certification, and nothing here is
      safety-certified for mains, medical, automotive or aerospace use.
\end{itemize}

\hsprovenance{The numbers are Table 1 and the board files behind it. The order
path is the ordering subagent and its scripts in .claude/skills/hwde.}

\end{document}
